\documentclass[11pt]{amsart}
\usepackage[T1]{fontenc}
\usepackage[utf8]{inputenc}
\usepackage{amssymb}
\usepackage[margin=1.05in]{geometry}
\usepackage[hidelinks]{hyperref}
\newcommand{\stamptime}{\the\numexpr\time/60\relax:\ifnum\numexpr\time-60*(\time/60)\relax<10 0\fi\the\numexpr\time-60*(\time/60)\relax}
\AddToHook{shipout/foreground}{\put(0.5\paperwidth,-\paperheight+0.55cm){\makebox[0pt]{\tiny\ttfamily unit/doc/\jobname.tex \ --- \ compiled \today, \stamptime}}}
\newtheorem{theorem}{Theorem}
\newtheorem{proposition}[theorem]{Proposition}
\newtheorem{lemma}[theorem]{Lemma}
\newtheorem{corollary}[theorem]{Corollary}
\theoremstyle{definition}
\newtheorem{definition}[theorem]{Definition}
\theoremstyle{remark}
\newtheorem{remark}[theorem]{Remark}
\newcommand{\Om}{\Omega}
\newcommand{\Map}{\mathrm{Map}_*}
\newcommand{\Z}{\mathbb{Z}}
\newcommand{\code}[1]{\texttt{#1}}
\newcommand{\cT}{\mathsf{isTrunc}}
\newcommand{\cC}{\mathsf{isConn}}
\newcommand{\refl}{\mathsf{refl}}
\newcommand{\ap}{\mathsf{ap}}
\newcommand{\Hom}{\mathrm{Hom}}
\newcommand{\hc}[2]{[#1,#2]_*}
\newcommand{\Fib}{\mathrm{Fib}}
\newcommand{\Const}{\mathsf{Const}}
\newcommand{\tr}[1]{\|#1\|}
\title{$\hc{\Fib\kappa}{K(B,3)}=0$}
\date{28 October 2026}
\begin{document}
\maketitle

\noindent\textbf{Purpose.} Informal proof for \S3b of
\code{doc/pi3-s2-postnikov-handoff-6.txt}: for $P=\Fib\kappa$ and every
abelian group $B$, the set $\hc P{K(B,3)}$ is a point. Written before
the formal proof; numbered statements are meant to be declarations.
Foundations as in the development (univalence, images; function
extensionality derived); no new axiom; the sphere axiom is not used.
\emph{Refereed} independently twice (a fresh reader given the note and the
library, not the conclusion; and a second reader): correct, with fixes
now incorporated (the name \code{STr.Q.relates}; the orientation of
Lemma~\ref{l:J}(a) and of the homotopy in Proposition~\ref{p:trans}, a
\code{sym} each; the endpoint generalisation in Lemma~\ref{l:J}(b); the
truncation index in Lemma~\ref{l:climb}, passed explicitly; the pointing
$\zeta$ in Proposition~\ref{p:exact} given directly; a shared spreading
lemma for Steps~1--3 of Propositions~\ref{p:trans} and~\ref{p:exact}).
\emph{Status: checked} (\code{sag/pi3-s2-postnikov/FibreH3.seg}, \S\ref{sec:corr}).

\smallskip\noindent\textbf{Hypotheses.} $\kappa:K(\Z,2)\to_*K(\Z,4)$ is
any pointed map for which $\Phi_B:\Hom(\Z,B)\to\hc{K(\Z,2)}{K(B,4)}$,
$g\mapsto[\kappa;K(g,4)]$ is injective for every $B$ (we use only the
injectivity in (H4), \code{KZ2Cohomology.H4}); it is instantiated with
\code{CupSquare.kappa} and \code{H4Proof.h4}. (H3)
(\code{H3Proof.h3}) is used for the base.

\smallskip\noindent\textbf{The idea.} Classically: the Serre sequence
$0=H^3(K(\Z,2))\to H^3(P)\to H^3(K(\Z,3))\xrightarrow{\tau}H^4(K(\Z,2))$,
$\tau=\Phi_B$. We prove the two facts that are needed, for this fibration
only, with explicit null-homotopies: (i) if the restriction of
$F:P\to_*K(B,3)$ to the fibre over the base point is $K(g,3)$, then
$\kappa;K(g,4)$ is null, by a homotopy built from $F$
(Proposition~\ref{p:trans}: this is ``$\tau(\text{restriction of }F)=0$'',
and it identifies the connecting map with $\Phi_B$ \emph{by
construction}: no separate connecting map is defined); (ii) if the
restriction is null, $F$ factors through $\mathrm{pr}:P\to K(\Z,2)$
(Proposition~\ref{p:exact}). Exactness of the sequence at
$H^3(K(\Z,3))$ is not needed.

\section{Conventions}

As in \code{doc/h4.pdf}: \code{isTrunc n X} is ``$X$ is
$(n-2)$-truncated'', \code{isConn n X} is ``$X$ is $(n-2)$-connected'';
$K_m=K(\Z,m)$, $K^B_m=K(B,m)$ (\code{EMMaps.KK}), base points $*$;
$p\cdot q$ is \code{trans p q} (by $J$ on $q$, so $p\cdot\refl\equiv p$);
$\Om Y=(y_0=y_0,\refl)$; $F;G$ is \code{ptcomp F G} ($F$ first); $0$ is
\code{ptconst}; $\Om F$ is \code{Maps.loop F} with
$(\Om F)(q)=F_0^{-1}\cdot(\ap_F q\cdot F_0)$ (\code{loop-ap}).
$e_m:\Om K^B_{m+1}\simeq_*K^B_m$ (\code{KK.e}).

Fix $B$, and write $W=K^B_4$, $w_0=*$. Then
\[P=\Fib\kappa=(x:K_2)\times T_x,\quad T_x=(\kappa x=*),\qquad
p_0=(*,\kappa_0),\]
with $\kappa_0:\kappa\,*=*$ the pointing of $\kappa$.
Levels: $K_2$ is $\cC\,3$; $K_4$ is $\cC\,5$, so each $T_x$ is $\cC\,4$
(\code{conn-id}); $W$ is $\cT\,6$, so $\Om W$ and each
$(G(\kappa x)=w_0)$ are $\cT\,5$; $W$ is $\cC\,5$, so $\Om W$ is $\cC\,4$
and $\Om^2W$ is $\cC\,3$ (both connected).

We work with targets $\Om W$ instead of $K^B_3$; the two are pointed
equivalent by $e_3$, and statements about classes transfer along a
pointed equivalence by pointed-equivalence induction (\code{ptEqInd}).

\section{Three small lemmas}

\begin{lemma}[constancy is a property]\label{l:const}
Let $T$ be $\cC\,4$, $Z$ be $\cT\,5$ and $d:T\to Z$. Then
\[\Const(d)=(z:Z)\times(\lambda\_.\,z=d)\]
is a proposition.
\end{lemma}
\begin{proof}
$\Const(d)$ is the fibre over $d$ of $\mathrm{c}:Z\to(T\to Z)$,
$z\mapsto\lambda\_.z$; so it suffices that $\mathrm{c}$ is an embedding
(\code{Groupoid.emb-fib-prop}). For $z,z':Z$,
$\mathsf{happly}\circ\ap_{\mathrm{c}}$ is, by $J$, the constant map
$(z=z')\to(T\to(z=z'))$, an equivalence by \code{ConnectedMaps.null}
($T$ is $\cC\,4$, $z=z'$ is $\cT\,4$); $\mathsf{happly}$ is an
equivalence; so $\ap_{\mathrm{c}}$ is one (\code{Whitehead.cancel-equiv}).
\end{proof}
Here the index matters: $T$ is 2-connected and $Z$ is 3-truncated, so
maps $T\to Z$ need not be constant; only the \emph{defect} of a map lies
in the 2-truncated $z=z'$, and that is what is used.

\begin{lemma}[pointing irrelevance, mere]\label{l:pi}
If $\Om(Z,z_0)$ is connected, $k,k':X\to Z$, $e:k=k'$, $\zeta:k\,x_0=z_0$,
$\zeta':k'x_0=z_0$, then $\tr{(k,\zeta)=(k',\zeta')}$.
\end{lemma}
\begin{proof}
$J$ on $e$; then $\zeta'=\zeta\cdot(\zeta^{-1}\cdot\zeta')$ and
$\zeta^{-1}\cdot\zeta'$ is merely $\refl$ (as \code{H4Classes.pi-set},
without the set hypothesis, since the goal is a mere proposition).
\end{proof}

\begin{lemma}[two $J$ computations]\label{l:J}
Let $Y$ be a type, $y,y_0:Y$, $k:y=y_0$.
\begin{enumerate}
\item[(a)] If $h:(y=y_0)\to C$, $c:C$ and $h(k\cdot q)=c$ for all
  $q:y_0=y_0$, then $h\,t=c$ for all $t:y=y_0$.
\item[(b)] Let $G:(Y,y_0)\to_*(W',w')$ and $h:(y=y_0)\to(w'=w')$ with
  $h(k\cdot q)=(\Om G)(q)$ for all $q$. Then for all $t:y=y_0$,
  \[\ap_Gk\cdot G_0\ =\ (\ap_Gt\cdot G_0)\cdot(h\,t)^{-1}.\]
\end{enumerate}
\end{lemma}
\begin{proof}
$J$ on $k$ ($y$ free; $h$ and the hypothesis in the motive). At
$k=\refl$: (a) $h\,t=h(\refl\cdot t)=c$ (\code{trans-refl-l}).
(b) With $u=\ap_Gt\cdot G_0$: $h\,t=h(\refl\cdot t)=(\Om G)(t)=G_0^{-1}\cdot u$.
The groupoid identity used is: for $u,v:c=d$,
\[u\cdot(v^{-1}\cdot u)^{-1}=v,\]
by $J$ on $u$ with the endpoint $d$ free and $v$ generalised in the motive
(at $d\equiv c$, $u=\refl$: $\refl\cdot(v^{-1}\cdot\refl)^{-1}\equiv\refl\cdot(v^{-1})^{-1}
=(v^{-1})^{-1}=v$, \code{trans-refl-l}, \code{sym-sym}). With $v=G_0$ the
right side is $(\ap_Gt\cdot G_0)\cdot(h\,t)^{-1}$, and the left side of the
claim is $\ap_G\refl\cdot G_0\equiv\refl\cdot G_0=G_0$ (\code{trans-refl-l}).
\end{proof}

\section{The fibre restriction}

\begin{definition}\label{d:res}
For $F:P\to_*\Om W$, its restriction to the fibre over $*$ is
\[\rho(F)=(q\mapsto F(*,\kappa_0\cdot q),\ F_0)\ :\ \Om K_4\to_*\Om W,\]
an element of $\code{PMO}=\Map(\Om K_4,\Om W)$. (The pointing is $F_0$ by
definition: $\kappa_0\cdot\refl\equiv\kappa_0$.)
\end{definition}

\begin{lemma}[the delooping of the restriction]\label{l:deloop}
For each $F$ there is a unique $G_F:K_4\to_*W$ with $\rho(F)=\Om G_F$.
\end{lemma}
\begin{proof}
\code{DeloopMaps} Corollary~11: $\code{Unloop}\,\rho(F)$ is contractible, with
$c=2$ ($K_4$ is $\cC\,4$, by \code{conn-down}) and $W$ $\cT\,6=\cT(3+3)$.
$G_F$ and the path are the components of its centre.
\end{proof}

\section{The transgression of the restriction is zero}

\begin{proposition}[transgression]\label{p:trans}
Let $F:P\to_*\Om W$ and $G:K_4\to_*W$ with $\rho(F)=\Om G$. Then
$\tr{0=\kappa;G}$ in $\Map(K_2,W)$; hence $[\kappa;G]=[0]$.
\end{proposition}
\begin{proof}
For $x:K_2$ let $Z_x=(G(\kappa x)=w_0)$ and
\[d_x:T_x\to Z_x,\qquad d_x(t)=(\ap_Gt\cdot G_0)\cdot F(x,t)^{-1}.\]
\emph{Step 1.} $\Const(d_x)$ is a proposition (Lemma~\ref{l:const}: $T_x$
is $\cC\,4$, $Z_x$ is $\cT\,5$).
\emph{Step 2.} $\Const(d_*)$ holds, with $z=\ap_G\kappa_0\cdot G_0$:
apply Lemma~\ref{l:J}(b) with $y=\kappa\,*$, $k=\kappa_0$,
$h=F(*,-)$; its hypothesis $F(*,\kappa_0\cdot q)=(\Om G)(q)$ is
$\mathsf{happly}$ of $\rho(F)=\Om G$ (first components; by
Definition~\ref{d:res}, $\rho(F)(q)$ is $F(*,\kappa_0\cdot q)$); then
function extensionality.
\emph{Step 3.} By \code{conn-elim} ($K_2$ is 0-connected, Step~1, Step~2),
$\Const(d_x)$ for every $x$. Let $N_x:G(\kappa x)=w_0$ be its first
component.
\emph{Step 4.} $N$ is a homotopy from the underlying map $x\mapsto
G(\kappa x)$ of $\kappa;G$ to $\lambda\_.w_0$; by Lemma~\ref{l:pi}
($\Om W$ is connected) $\tr{0=\kappa;G}$, and
\code{STr.relates} gives $[0]=[\kappa;G]$.
\end{proof}

Note that only the first component $N_x$ is used: the pointing of the
null-homotopy is irrelevant by Lemma~\ref{l:pi}. What is used about $F$
is exactly that $d_x$ is constant on each fibre: the constancy is a
proposition, holds over $*$ by the hypothesis on $\rho(F)$, and spreads
along the connected base. This is the transgression: $d_x$ compares the
fibre map $F(x,-)$ with the canonical map $t\mapsto\ap_Gt\cdot G_0$ from
$T_x$ to the fibre $Z_x$ of $W$, and its constant value trivialises
$Z_x$.

\section{Classes with null restriction come from the base}

\begin{proposition}[exactness at $H^3(P)$]\label{p:exact}
If $F:P\to_*\Om W$ and $\rho(F)=0$, then $\tr{0=F}$.
\end{proposition}
\begin{proof}
\emph{Step 1.} For each $x$, $\Const(F(x,-))$, with $F(x,-):T_x\to\Om W$,
is a proposition (Lemma~\ref{l:const}; $\Om W$ is $\cT\,5$).
\emph{Step 2.} $\Const(F(*,-))$ holds with $z=\refl$:
Lemma~\ref{l:J}(a) with $k=\kappa_0$, $c=\refl$, the hypothesis being
$\mathsf{happly}$ of $\rho(F)=0$.
\emph{Step 3.} By \code{conn-elim}, $\Const(F(x,-))$ for all $x$; let
$H\,x:\Om W$ be its first component, so $H\,x=F(x,t)$ for all $t$.
\emph{Step 4.} The pointing is given directly:
$\zeta=(H\,*=F(*,\kappa_0))\cdot F_0:H\,*=\refl$ (the constancy witness at
$t=\kappa_0$, then the pointing of $F$, since $p_0=(*,\kappa_0)$);
$(H,\zeta):K_2\to_*\Om W$. By (H3), transferred along
$e_3^{-1}:K^B_3\simeq_*\Om W$, $\hc{K_2}{\Om W}$ is a point, so merely
$e:0=(H,\zeta)$ (\code{effective}).
\emph{Step 5.} From $e$, $\refl=H\,x$ for all $x$, so
$\refl=F(x,t)$ for all $(x,t)$, i.e.\ $\lambda\_.\refl=F_1$; by
Lemma~\ref{l:pi} ($\Om^2W$ is connected), $\tr{0=F}$.
\end{proof}

\section{Coefficients: every $K_4\to_*W$ is some $K(g,4)$}

\begin{lemma}[climbing]\label{l:climb}
For $m\ge1$ and $G:K_{m+1}\to_*K^B_{m+1}$ let
$G^\flat=e_m^{-1};\Om G;e_m:K_m\to_*K^B_m$ (\code{unconj}). Then
$\code{conj}_m(G^\flat)=\Om G$, where $\code{conj}_m(F)=e_m;F;e_m^{-1}$
(\code{Lev.conj}). If $G^\flat=K(g,m)$, then $K(g,m+1)=G$.
\end{lemma}
\begin{proof}
The first equation: associativity, \code{pt-inv-l} twice, the unit
laws (as \code{H4Bilinear.H2.conj-unconj}, which is the case $m=1$). Then
$\Om K(g,m+1)=\code{conj}_m(K(g,m))$ (\code{Lev.loopK})
$=\code{conj}_m(G^\flat)=\Om G$, and \code{unique-loop} at some $c$ with
$K_{m+1}$ $\cC(c+2)$ and $K^B_{m+1}$ $\cT(2c+2)$: for $m=2$, $c=2$ and
$\cT\,5\Rightarrow\cT\,6$ (\code{trunc-suc}); for $m=3$, $c=2$ with
\code{conn-down} ($\cC\,5\Rightarrow\cC\,4$) and $\cT\,6$ exactly. The
lift is a step: $\cT(m+3)$ is not literally $\cT(2c+2)$.
\end{proof}

\begin{corollary}\label{c:surj4}
Every $G:K_4\to_*K^B_4$ is $K(g,4)$ for $g=\mathrm{inv}_4(G)$, where
$\mathrm{inv}_4(G)=\mathrm{inv}_2(((G^\flat)^\flat))$ and $\mathrm{inv}_2$
is \code{H4Bilinear.H2.inv2} (\code{k2-surj}: $K(\mathrm{inv}_2F,2)=F$).
\end{corollary}
\begin{proof}
Lemma~\ref{l:climb} at $m=2$, then at $m=3$.
\end{proof}

\section{The theorem}

\begin{theorem}\label{t:main}
For every abelian $B$, $\hc P{K^B_3}$ is contractible, with centre $[0]$.
\end{theorem}
\begin{proof}
By transport along $e_3:\Om W\simeq_*K^B_3$ it suffices to treat
$\hc P{\Om W}$. Classes are eliminated into the propositions $[0]=c$
(\code{STr.ind}); for $F:P\to_*\Om W$ we show $\tr{0=F}$.
Let $G=G_F$ (Lemma~\ref{l:deloop}). By Proposition~\ref{p:trans},
$[\kappa;G]=[0]$. By Corollary~\ref{c:surj4}, $G=K(g,4)$ and
$0=K(g',4)$ with $g=\mathrm{inv}_4G$, $g'=\mathrm{inv}_40$. Then
\[\Phi_B(g)=[\kappa;K(g,4)]=[\kappa;G]=[0]=[\kappa;0]=[\kappa;K(g',4)]=\Phi_B(g')\]
(\code{Lev.K-def} for the first and last steps, as in
\code{H4Proof.phi-diag}; \code{comp-const-r} for $[0]=[\kappa;0]$). By
injectivity of $\Phi_B$, $g=g'$, so $G=K(g,4)=K(g',4)=0$ and
$\rho(F)=\Om G=\Om0=0$ (\code{loop-const}). Proposition~\ref{p:exact}
gives $\tr{0=F}$.
\end{proof}

\begin{remark}
(1) The connecting map is never defined separately: the step
``$\tau(\rho F)=\Phi_B(g)$'' is Proposition~\ref{p:trans} applied to
$G=K(g,4)$, and the class it produces is literally
$[\kappa;K(g,4)]$, the value of \code{KZ2Cohomology.Phi}.
(2) Only injectivity of $\Phi_B$ is used, and only (H3) for $K(\Z,2)$; no
statement about $H^3(K(\Z,3);B)$ beyond Lemma~\ref{l:deloop} and
Corollary~\ref{c:surj4}. (3) Nothing is said about the type of
null-homotopies; all conclusions are mere.
(4) The same proof works for any $\kappa$ with $\Phi_B$ injective; the
module is parameterised by $\kappa$ and a proof of (H4).
\end{remark}

\section{Correspondence with the formal proof}\label{sec:corr}
\emph{Status: checked} (69 modules, 0 failed; a fresh process: 2877
declarations, 1.1\,s; every \code{FibreH3} declaration under 0.01\,s;
\code{fibre.opens}, 84 entries, every line ``nothing else opened'').
Module \code{FibreH3}: \code{Const}, \code{const-emb}, \code{const-prop}
(Lemma~\ref{l:const}); \code{pi-mere} (Lemma~\ref{l:pi});
\code{reparam}, \code{cancel-v}, \code{TG.tg} (Lemma~\ref{l:J});
\code{Climb[A,B]}: \code{unconj}, \code{conj-unconj}, \code{climb},
\code{inv3}, \code{surj3}, \code{inv4}, \code{surj4}
(Lemma~\ref{l:climb}, Corollary~\ref{c:surj4}); \code{spread} (Steps~1--3
of Propositions~\ref{p:trans}, \ref{p:exact});
\code{FibAt[kappa, h4k, B]}: \code{T}, \code{cT}, \code{cOW}, \code{cOOW}
(levels), \code{res}, \code{res-def} (Definition~\ref{d:res}),
\code{deloop-c}, \code{deloop}, \code{deloop-eq} (Lemma~\ref{l:deloop}),
\code{Trans[F,G,w]}: \code{Z}, \code{d}, \code{hyp}, \code{const-base},
\code{Nu} ($N_x$), \code{null}, \code{trans-cls}
(Proposition~\ref{p:trans}); \code{h3-loop} ((H3) transferred);
\code{Exact[F,z]}: \code{hyp}, \code{const-base}, \code{cF}, \code{H},
\code{H-eq}, \code{zeta}, \code{base-zero}, \code{from-base}, \code{exact}
(Proposition~\ref{p:exact}); \code{g0}, \code{Kill[F]}: \code{g},
\code{cls-eq}, \code{g-eq}, \code{G-zero}, \code{res-zero},
\code{merely-zero}, \code{contr-loop}, \code{contr}
(Theorem~\ref{t:main}); \code{fib-h3 : (B : Ab) -> isContr [Fib kappa, K(B,3)]}
for \code{CupSquare.kappa} and \code{H4Proof.h4}.
\end{document}
